%=====================================================================
\section{Reinforcement learning in a nutshell}\label{sec:primer}
%=====================================================================

This section explains the ideas behind reinforcement learning (RL) for a
reader who knows hydrologic modeling and calibration but has not used RL.
Everything later in the report builds on it. The glossary at the end of the
report collects the terms.

\subsection{What problem does RL solve?}
Think of a reservoir operator. Every morning they look at the reservoir
level, the current inflow and the weather forecast, and decide how much
water to release today. Tomorrow they see the consequences and decide
again. The goal is not to make each single day look good. It is to do well
over the whole season: to meet demand, avoid floods and not waste water
over the spillway. The operator does not know the future storms, so a good
decision now must hedge against many possible futures.

This is a \emph{sequential decision problem under uncertainty}. RL is a
family of methods that \emph{learn a decision rule for such problems by
trial and error}: they try decisions in a simulator, observe what happens
and gradually improve the rule. In this report, the simulator is an
OpenHydroQual model.

\paragraph{How this differs from calibration.} In model calibration we
search for \emph{parameter values} that make a model fit observed data
(for example with a genetic algorithm or Levenberg--Marquardt). In RL we
search for a \emph{decision rule}, called a \emph{policy}: a function that
maps what the operator sees to what the operator does. The rule is judged
by how well it performs over many simulated futures, not by how well it
fits data.

\paragraph{Why not just optimize the release schedule?} One could fix a
storm sequence and optimize the 90 daily releases for it. But real storms
are unknown in advance. A schedule that is optimal for one storm sequence
can be disastrous for another: it might release water ahead of a storm that
never comes. What we need is a rule that \emph{reacts} to what it observes
(closed-loop control) instead of a fixed plan (open-loop control). Classical
alternatives each have a catch:
\begin{itemize}\itemsep1pt
\item \textbf{Rule curves} are hand-designed rules with a few parameters
(for example, ``release extra water above 80\% full''). They are simple and
transparent but limited to the shape the designer chose.
\item \textbf{Stochastic dynamic programming} finds the optimal rule
exactly, but its cost explodes with the number of state variables (the
``curse of dimensionality''). The cascade in Section~\ref{sec:cascade} has
16 observed quantities, far beyond what it can handle.
\item \textbf{Model predictive control} re-optimizes a plan every day over
a forecast horizon. It needs a fast optimizer inside the loop and a model
of forecast uncertainty.
\end{itemize}
RL instead learns a flexible rule (here a small neural network) directly
from simulations. The price is that it needs a lot of them, hundreds of
thousands of simulated days here, which is why simulator speed matters
(Section~\ref{sec:performance}).

\subsection{The agent--environment loop}
RL describes the problem as a loop between an \emph{agent} (the decision
maker) and an \emph{environment} (everything else) (Figure~\ref{fig:loop}).
At each step $t$:
\begin{enumerate}\itemsep1pt
\item the agent receives an \textbf{observation} $o_t$, the information
available for the decision (storage, inflow, forecasts, demand, \dots);
\item it chooses an \textbf{action} $a_t$ (here, today's release);
\item the environment advances one step (here, OHQ simulates one day) and
returns a \textbf{reward} $r_t$, one number that scores how good that step
was, together with the next observation $o_{t+1}$.
\end{enumerate}
A sequence of steps from a start to an end is an \textbf{episode}; here one
episode is one 90-day season with its own random storms. The sum of the
rewards over an episode is the \textbf{return}. The agent's goal is to
choose actions that make the return as large as possible \emph{on average
over all the seasons it might face}.

\begin{figure}[h]
\centering
\begin{tikzpicture}[>=Stealth, font=\small,
  box/.style={draw, rounded corners, align=center, minimum height=1.5cm, text width=3.4cm, fill=#1}]
  \node[box=blue!8] (agent) at (0,0) {\textbf{Agent} (policy $\pi$)\\the operator's decision rule: a neural network};
  \node[box=green!8] (env) at (9.6,0) {\textbf{Environment}\\OHQ reservoir model + random storms and demand};
  \draw[->, thick] ([yshift=4mm]agent.east) -- node[above, align=center, font=\footnotesize]{action $a_t$ (release today)} ([yshift=4mm]env.west);
  \draw[->, thick] ([yshift=-4mm]env.west) -- node[below, align=center, font=\footnotesize]{reward $r_t$ (score of the day)\\observation $o_{t+1}$ (what is seen next)} ([yshift=-4mm]agent.east);
\end{tikzpicture}
\caption{The agent--environment loop, repeated once per simulated day.}
\label{fig:loop}
\end{figure}

\paragraph{Designing the reward is designing the objective.} The agent
optimizes exactly what the reward measures and nothing else. Here the
reward penalizes supply deficits, downstream flooding, spill and ending
the season empty (Section~\ref{sec:objective}). If a penalty were missing
the agent would ignore that concern, and if a weight were too large it
would over-react to it. Choosing the reward is the main place where
engineering judgment enters.

\paragraph{Observation versus state.} The \emph{state} is everything that
determines the future. The \emph{observation} is what the agent actually
sees. The agent never sees future storms, only noisy forecasts of them, so
it must act well despite incomplete information, exactly like a real
operator.

\paragraph{Discounting.} Many RL methods maximize a \emph{discounted}
return, $r_0+\gamma r_1+\gamma^2 r_2+\dots$, with a discount factor $\gamma$
slightly below 1 (here $0.99$). Rewards far in the future count a little
less. This makes learning more stable, because distant consequences are
hard to attribute to today's action. The price is a mild preference for
the near term (discussed in Section~\ref{sec:objective}).

\subsection{The policy}
The \textbf{policy} $\pi$ is the decision rule: given $o_t$, it returns
$a_t$. A rule curve is a policy with one or two parameters. RL uses a much
more flexible function, a small \textbf{neural network}. Here it is a stack
of two layers of 64 simple nonlinear units, which can represent smooth
functions of all the observed quantities at once. Its adjustable numbers,
the \emph{weights}, collectively written $\theta$, play the role that model
parameters play in calibration. Training an RL agent means adjusting
$\theta$.

During training the policy is \textbf{stochastic}: the network outputs a
\emph{mean} release, and the actual release is drawn from a bell curve
around it. This random perturbation is how the agent \textbf{explores},
trying slightly different actions to discover which work better. The
width of the bell curve, $\sigma$, is learned too, and it shrinks as the
agent becomes confident (see the learning curves). When we \emph{use} the
trained policy for evaluation, we switch the randomness off and apply the
mean action.

\subsection{How learning works, in one paragraph}
The agent plays many seasons in the simulator with its current policy and
records what happened. For every decision it asks: \emph{did things turn
out better or worse than expected after this action?} It then adjusts
$\theta$ so that actions that turned out better than expected become more
likely, and actions that turned out worse become less likely. Then it plays
more seasons with the improved policy, and so on. Two difficulties make
this harder than it sounds:
\begin{itemize}\itemsep1pt
\item \textbf{Credit assignment.} A flood on day 50 may be the fault of a
decision on day 40 not to release water. The agent must work out which
earlier actions deserve the blame. It does this with a \emph{value
function} (next paragraph).
\item \textbf{Noise.} Two seasons can differ wildly in storminess, so a
bad outcome may be bad luck rather than a bad action. Comparing against
``what was expected'' filters out much of this luck.
\end{itemize}

\paragraph{The value function (the ``critic'').} Alongside the policy, a
second neural network $V_\phi(o)$ learns to predict the return still to
come from a given situation: roughly, ``how good is it to be here?''. It
is like a forecast of the remaining season's score given today's storage,
inflows and forecasts. Comparing what actually happened after an action
with this prediction tells the agent whether the action was better or
worse than expected. The policy network is then called the \emph{actor}
and the value network the \emph{critic}, and the method an
\emph{actor--critic} method. Section~\ref{sec:ppo} gives the details for
the specific algorithm used, PPO.

\subsection{The same ideas in equations}\label{sec:primer_math}
This subsection states the ideas above precisely. It assumes comfort with
probability and calculus, and makes the link to stochastic dynamic
programming (SDP), which many water resources engineers know from
reservoir operation.

\paragraph{Markov decision process.} The problem is formalized as a
\emph{Markov decision process} (MDP): states $s\in\mathcal S$, actions
$a\in\mathcal A$, a transition law $p(s_{t+1}\mid s_t,a_t)$, a reward
function $r_t=r(s_t,a_t,s_{t+1})$ and a horizon $T$. ``Markov'' means the
next state depends only on the current state and action, not on the
earlier history. In our problem the state contains the storage, the
current and future storm hydrographs and the demand. The transition law is
the OHQ mass balance plus the random weather. Crucially, we \emph{never
write $p$ down}: we can only \emph{sample} from it by running the
simulator. This is the key practical difference from SDP, which needs $p$
explicitly as transition probabilities on a discretized grid.

\paragraph{Policy and objective.} A stochastic policy $\pi_\theta(a\mid o)$
is a probability density over actions given the observation, with
parameters $\theta$. In this report it is a normal density,
\[
\pi_\theta(a\mid o)=\frac{1}{\sqrt{2\pi}\,\sigma}\exp\!\Big(-\frac{(a-\mu_\theta(o))^2}{2\sigma^2}\Big),
\]
whose mean $\mu_\theta(o)$ is the output of the neural network. A
\emph{trajectory} $\tau=(o_0,a_0,r_0,o_1,a_1,r_1,\dots)$ is one simulated
season, and its discounted return is $G(\tau)=\sum_{t=0}^{T-1}\gamma^t r_t$.
The goal is to find
\begin{equation}
\theta^\star=\arg\max_\theta\; J(\theta),\qquad J(\theta)=\mathbb E_{\tau\sim\pi_\theta}\big[G(\tau)\big],
\label{eq:rl_objective}
\end{equation}
where the expectation is over the random weather \emph{and} the policy's
own randomness. This is an optimization over $\theta$, just like
calibration. The difference is that the objective is an expectation we can
only estimate by simulating seasons.

\paragraph{Value functions and the Bellman equation.} The \emph{value} of a
state under policy $\pi$ is the expected return still to come,
\[
V^\pi(s_t)=\mathbb E_\pi\Big[\sum_{k=0}^{T-1-t}\gamma^k r_{t+k}\;\Big|\;s_t\Big],
\]
and the \emph{action value} $Q^\pi(s,a)$ is the same quantity when action
$a$ is taken first and $\pi$ is followed afterwards. Splitting off the
first reward gives the \emph{Bellman equation}
\begin{equation}
V^\pi(s_t)=\mathbb E\big[r_t+\gamma\,V^\pi(s_{t+1})\big],
\qquad
V^\star(s_t)=\max_{a}\;\mathbb E\big[r_t+\gamma\,V^\star(s_{t+1})\;\big|\;s_t,a\big].
\label{eq:bellman}
\end{equation}
The right-hand equation is exactly the backward recursion of SDP for
reservoir operation: the best value today is the best immediate reward plus
the expected (discounted) value of tomorrow's storage. SDP solves it by
discretizing the state (for example storage into 50 classes and inflow into
10) and sweeping backward in time. This is exact but needs the transition
probabilities, and the grid grows exponentially with the number of state
variables. With the 16 observed quantities of the cascade problem, a grid
of only 10 levels each would have $10^{16}$ cells. RL keeps the Bellman
logic but (i) replaces the grid by neural networks that approximate $V$ and
$\pi$ smoothly, and (ii) replaces the transition probabilities by simulated
samples. In this sense RL is \emph{approximate, simulation-based dynamic
programming}.

\paragraph{The policy gradient.} To improve $\theta$ by gradient ascent we
need $\nabla_\theta J$, but $J$ is an expectation over simulator runs, and
the simulator (OHQ) is not differentiable. The \emph{log-derivative trick}
gets around this. Write the probability of a trajectory as
$P_\theta(\tau)=p(s_0)\prod_t\pi_\theta(a_t\mid o_t)\,p(s_{t+1}\mid s_t,a_t)$.
Then
\[
\nabla_\theta J=\nabla_\theta\!\int\!P_\theta(\tau)G(\tau)\,d\tau
=\int\!P_\theta(\tau)\,\nabla_\theta\log P_\theta(\tau)\,G(\tau)\,d\tau
=\mathbb E_{\tau}\Big[G(\tau)\sum_t\nabla_\theta\log\pi_\theta(a_t\mid o_t)\Big],
\]
because $\nabla P=P\,\nabla\log P$, and the transition terms $p(\cdot)$ do
not depend on $\theta$, so their log-derivatives vanish. \emph{The physics
drops out}: we only need to differentiate our own policy network, never the
simulator. The expectation can be estimated by averaging over simulated
seasons. In words: push up the log-probability of every action in a season
in proportion to how good that season was.

\paragraph{Reducing the noise: baselines and advantages.} That estimator is
correct but extremely noisy: a stormy season punishes every action in it,
good or bad. Two refinements fix most of this. (i) An action cannot affect
rewards that came before it, so only the reward \emph{after} step $t$ is
used to judge $a_t$. (ii) Subtracting any function of the observation,
$b(o_t)$, does not change the expected gradient, because
$\mathbb E_{a\sim\pi}[\nabla_\theta\log\pi_\theta(a\mid o)]=\nabla_\theta\!\int\!\pi_\theta\,da=\nabla_\theta 1=0$.
But a good choice of $b$ greatly reduces the variance. The best simple
choice is the value function itself, which gives
\begin{equation}
\nabla_\theta J=\mathbb E\Big[\sum_t \nabla_\theta\log\pi_\theta(a_t\mid o_t)\;A^\pi(o_t,a_t)\Big],
\qquad A^\pi(o_t,a_t)=Q^\pi(o_t,a_t)-V^\pi(o_t).
\label{eq:pg}
\end{equation}
The \textbf{advantage} $A$ measures how much better action $a_t$ was than
what the policy does on average in that situation. This is the precise
form of ``did things turn out better or worse than expected?''. Positive
advantage: make that action more likely. Negative: less likely.

\paragraph{Estimating the advantage: the TD error.} We do not know $Q^\pi$
or $V^\pi$. We learn an approximation $V_\phi\approx V^\pi$ (the critic)
by least-squares regression on observed returns. Then the one-step Bellman
relation~\eqref{eq:bellman} gives a cheap advantage estimate, the
\emph{temporal-difference (TD) error}
\begin{equation}
\delta_t=r_t+\gamma V_\phi(o_{t+1})-V_\phi(o_t).
\label{eq:td}
\end{equation}
Read it as ``what actually happened today, plus what the critic now
expects for the rest of the season, minus what the critic expected
yesterday''. If $\delta_t>0$, the day went better than expected. Averaging
it over many steps is how credit is assigned through time:
Section~\ref{sec:ppo} shows how PPO blends several steps of $\delta$ into a
better estimate.

\subsection{Training versus testing}
As in calibration and validation, an agent must be judged on data it was
not trained on. Each season's random storms are generated from a
\emph{seed} (a number that fixes the random sequence). Training uses one
set of seeds. Testing uses 30 different seeds, the same 30 for every
policy, so that all policies face identical storms and the comparison is
fair.

\subsection{What the agent sees in this report}
\begin{center}
\small
\begin{tabular}{p{0.24\textwidth}p{0.68\textwidth}}
\toprule
RL concept & In the single-reservoir problem\\
\midrule
environment & OHQ reservoir model + storm and demand generator\\
step & one day (one OHQ run of one day)\\
episode & one 90-day season\\
observation $o_t$ & storage, current inflow, 1--3-day inflow forecasts, demand, days left\\
action $a_t$ & release for the day, from 0 to 6{,}000 m$^3$/d\\
reward $r_t$ & minus the day's penalties (deficit, flooding, spill)\\
return & minus the season's total penalty (closer to 0 is better)\\
policy & neural network mapping $o_t$ to $a_t$\\
\bottomrule
\end{tabular}
\end{center}
